%% 
%% Template article for cas-sc documentclass for 
%% double column output.

\documentclass[a4paper,fleqn]{cas-sc}

% If the frontmatter runs over more than one page
% Use the longmktitle option.

%\documentclass[a4paper,fleqn,longmktitle]{cas-sc}

%\usepackage[numbers]{natbib}
%\usepackage[authoryear]{natbib}
\usepackage[numbers]{natbib}
%% The amssymb package provides various useful mathematical symbols
\usepackage{amssymb}
%% The amsmath package provides various useful equation environments.
\usepackage{amsmath}
\usepackage{epstopdf}
\usepackage{mathptmx,mathtools}
\usepackage{subcaption}
\usepackage{float} % Add this in your preamble
\usepackage[export]{adjustbox}
\usepackage{array} % Needed for m{} column type
\usepackage[percent]{overpic} % <-- put this in your preamble



\usepackage{amsmath}
\usepackage{nicefrac}
\usepackage{dirtytalk}
\usepackage{graphicx}   % For including images
\usepackage{tikz}
\usetikzlibrary{calc,tikzmark} % needed <<<<<<<<<<
\usepackage{xcolor}     % For coloring
\usepackage{lineno}
\renewcommand\linenumberfont{\normalfont\bfseries\small} % <-- Increase line number size
\usetikzlibrary{arrows.meta}
\usetikzlibrary{calc, positioning}
%%%Author macros
\def\tsc#1{\csdef{#1}{\textsc{\lowercase{#1}}\xspace}}
\tsc{WGM}
\tsc{QE}

\begin{document}
	
	\setcounter{figure}{1}
	
	\begin{figure}[!t]
	\centering
	
	% ------- Subfigure (a) ------- %
	\begin{subfigure}[t]{0.48\textwidth}
		\centering
		\begin{overpic}[width=\textwidth,trim={3cm 4cm 0cm 0cm},clip]{Fig2a.eps}
			% Panel label (edit coords as needed)
			\put(-5,85){\fontsize{15}{12}\selectfont\fontfamily{ptm}\selectfont\textbf{(a)}}
			
			\put(14,7){\fontfamily{ptm}{\textbf{Eulerian}}}
			\put(49,7){\fontfamily{ptm}{\textbf{Moving frame}}}
			% Example extra labels:
			% \put(20,80){\small PRA}
			% \put(60,20){\small Streamlines}
		\end{overpic}
	\end{subfigure}%
	% ------- Subfigure (b) ------- %
	\begin{subfigure}[t]{0.5\textwidth}
		\centering
		\begin{overpic}[width=\textwidth,trim={3cm 4.5cm 0cm 0cm},clip]{Fig2b.eps}
			\put(-5,82){\fontsize{15}{12}\selectfont\fontfamily{ptm}\selectfont\textbf{(b)}}
			\put(14,7){\fontfamily{ptm}{\textbf{Eulerian}}}
			\put(49,7){\fontfamily{ptm}{\textbf{Moving frame}}}
			% \put(15,75){\small Weber (1980)}
		\end{overpic}
	\end{subfigure}
	
	\caption{Primary wake of a rising bubble visualized using (a) streamlines and (b) PRA in two different frames of reference for bubble case $C_6$. The left side of the axisymmetric centerline corresponds to the Eulerian frame ($\mathbf{x}$), while the right side represents the Lagrangian frame ($\mathbf{y}$).}
	
	\label{Fig2}
	\end{figure}
%%%%%%%%%%%%%%%%%%%%%%%%%%%%%________Figure_3_end___________%%%%%%%%%%%%%%%%%%%%%
	
\end{document}